\chapter{Phân tích nghiệp vụ}

Chương này mô tả hệ thống ở góc nhìn nghiệp vụ: ai dùng, dùng để làm gì, quy trình chạy ra
sao và những quy tắc nào bắt buộc phải đúng. Toàn bộ thiết kế kỹ thuật ở các chương sau
đều truy nguyên về những nội dung trong chương này.

\section{Tác nhân}

\begin{longtable}{L{3.0cm} L{3.0cm} L{7.6cm}}
\toprule
\textbf{Tác nhân} & \textbf{Loại} & \textbf{Mục tiêu chính}\\
\midrule
\endhead
Người học & Chính, con người & Nâng bậc CEFR bằng việc luyện tập hằng ngày; nhận phản hồi đúng trình độ và đúng lỗi của mình.\\
Quản trị viên nội dung & Phụ, con người & Biên soạn, kiểm duyệt và xuất bản khoá học; bảo đảm mọi bài học đều chơi được.\\
Quản trị viên hệ thống & Phụ, con người & Quản lý tài khoản, vai trò, cấu hình mô hình AI, theo dõi sức khoẻ dịch vụ.\\
Trợ lý AI (Lexi) & Hệ thống & Chẩn đoán lỗi, sinh phản hồi có căn cứ, sinh quan sát học tập.\\
Tác tử nội dung & Hệ thống & Sinh khoá học và bài tập từ danh sách từ có nhãn CEFR.\\
Bộ lập lịch nền & Hệ thống & Nhắc ôn tập, tổng hợp thống kê, nạp trước nội dung, đặt lại giải đấu.\\
Nhà cung cấp LLM & Ngoài & Groq, Gemini, Ollama — cung cấp năng lực sinh ngôn ngữ.\\
Firebase & Ngoài & Xác thực mạng xã hội và thông báo đẩy.\\
Nguồn nội dung & Ngoài & RSS tin tức, YouTube, podcast, kho sách, kho ngữ liệu mở.\\
RevenueCat & Ngoài & Xử lý mua trong ứng dụng và quyền lợi gói trả phí.\\
\bottomrule
\caption{Tác nhân của hệ thống}
\end{longtable}

\section{Sơ đồ use case nghiệp vụ}

\begin{figure}[htbp]
\centering
\begin{tikzpicture}[scale=0.95, transform shape]
  \node[lane, minimum width=8.2cm, minimum height=8.6cm] (sys) at (0,0) {};
  \node[font=\scriptsize\bfseries, anchor=north, yshift=-3pt] at (sys.north) {Hệ thống LexiLingo};

  \node[uc] (u1) at (0, 3.2)  {Học theo lộ trình CEFR};
  \node[uc] (u2) at (0, 2.3)  {Ôn tập từ vựng theo lịch};
  \node[uc] (u3) at (0, 1.4)  {Hội thoại với trợ lý AI};
  \node[uc] (u4) at (0, 0.5)  {Luyện nói và phát âm};
  \node[uc] (u5) at (0,-0.4)  {Chơi trò chơi ngôn ngữ};
  \node[uc] (u6) at (0,-1.3)  {Học qua nội dung thực tế};
  \node[uc] (u7) at (0,-2.2)  {Theo dõi tiến độ, thành tích};
  \node[uc] (u8) at (0,-3.1)  {Đánh giá và thăng bậc};

  \node[uc, fill=cSoftO, draw=cGate] (a1) at (0,-4.0) {Biên soạn \& kiểm duyệt nội dung};

  \stickman{-5.6}{1.2}{Người học}
  \stickman{-5.6}{-2.4}{Quản trị viên\\nội dung}
  \stickman{5.6}{2.2}{Trợ lý AI\\(Lexi)}
  \stickman{5.6}{-1.0}{Nhà cung cấp\\LLM}
  \stickman{5.6}{-3.4}{Firebase /\\Nguồn nội dung}

  \foreach \t in {u1,u2,u3,u4,u5,u6,u7,u8}
    \draw[draw=cInfra!60, line width=0.5pt] (-5.35,1.2) -- (\t.west);
  \draw[draw=cGate!70, line width=0.5pt] (-5.35,-2.4) -- (a1.west);

  \draw[draw=cInfra!60, line width=0.5pt] (u3.east) -- (5.35,2.2);
  \draw[draw=cInfra!60, line width=0.5pt] (u4.east) -- (5.35,2.2);
  \draw[draw=cInfra!60, line width=0.5pt] (u8.east) -- (5.35,-1.0);
  \draw[draw=cInfra!60, line width=0.5pt] (a1.east) -- (5.35,-1.0);
  \draw[draw=cInfra!60, line width=0.5pt] (u6.east) -- (5.35,-3.4);
\end{tikzpicture}
\caption{Sơ đồ use case nghiệp vụ — tám nghiệp vụ người học và một nghiệp vụ quản trị nội dung}
\label{fig:uc-nghiepvu}
\end{figure}

\section{Chín nghiệp vụ chính}

\subsection{NV1 — Học theo lộ trình CEFR}

Người học ghi danh một khoá học, đi theo lộ trình dạng zigzag gồm các chương và bài học.
Mỗi bài học chứa một chuỗi bài tập; hoàn thành bài học tạo ra ba loại dữ liệu: bản ghi
hoàn thành, điểm XP và bằng chứng năng lực gửi vào hai động cơ mô hình hoá người học.

\textbf{Quy tắc nghiệp vụ ràng buộc:} một bài học chỉ mở được khi có ít nhất một bài tập
(mục~\ref{sec:qtnv}); khoá học chỉ xuất bản được khi mọi bài học đều thoả điều kiện đó.

\subsection{NV2 — Ôn tập từ vựng theo lịch}

Người học lưu từ vào bộ sưu tập (thủ công, hoặc lưu nhanh khi đang đọc nội dung). Hệ thống
tính hạn ôn cho từng khái niệm theo FSRS. Đến hạn, người học ôn bằng thẻ và tự chấm mức
nhớ 0--5; kết quả cập nhật độ bền, độ khó và mức nắm vững, từ đó sinh hạn ôn mới.

\subsection{NV3 — Hội thoại với trợ lý AI}

Người học nhắn tin hoặc nói với Lexi. Hệ thống chẩn đoán ý định và lỗi ngôn ngữ, truy xuất
bằng chứng từ đồ thị tri thức, sinh phản hồi có căn cứ theo chiến lược sư phạm phù hợp
trình độ, đồng thời ghi lại một quan sát học tập cho mỗi lượt.

\subsection{NV4 — Luyện nói và phát âm}

Người học nói vào micro; hệ thống nhận dạng theo dòng, phát hiện điểm kết thúc phát ngôn,
phân tích chuỗi âm vị, so với phát âm mục tiêu và trả về điểm cùng gợi ý sửa — có nhận
diện riêng các lỗi đặc trưng của người Việt.

\subsection{NV5 — Chơi trò chơi ngôn ngữ}

Sáu trò chơi mini rút dữ liệu từ ngân hàng từ vựng hoặc ngân hàng câu hỏi ngữ pháp. Kết
quả trò chơi được quy đổi sang kỹ năng CEFR tương ứng theo bảng ánh xạ cố định, không suy
đoán từ tên trò chơi.

\subsection{NV6 — Học qua nội dung thực tế}

Bốn nguồn: tin tức, podcast, sách phân cấp, video YouTube. Hệ thống nạp trước nội dung
theo lịch và sinh quiz kèm theo. \textbf{Chỉ phần quiz tạo ra bằng chứng học tập}; việc
nghe hoặc đọc thuần tuý cố ý không ghi nhận gì.

\subsection{NV7 — Theo dõi tiến độ và thành tích}

Người học xem chuỗi ngày học, bản đồ nhiệt hoạt động, điểm sáu kỹ năng, xu hướng 7 và
30 ngày, huy hiệu, bảng xếp hạng, ví và cửa hàng vật phẩm.

\subsection{NV8 — Đánh giá và thăng bậc}

Có hai đường vào một bậc CEFR: bài kiểm tra xếp lớp khi mới bắt đầu, và cổng thi thăng bậc
khi đã tích luỹ đủ bằng chứng. Cổng thi chặn việc lên quá một bậc mỗi lần.

\subsection{NV9 — Biên soạn và kiểm duyệt nội dung}

Quản trị viên tải lên danh sách từ có nhãn CEFR, chạy tác tử nội dung, rà soát kết quả
trong hàng đợi kiểm duyệt, sửa hoặc yêu cầu sinh lại, rồi mới xuất bản.

\section{Quy trình nghiệp vụ trọng tâm}

\subsection{Quy trình hoàn thành một bài học}

\begin{figure}[htbp]
\centering
\begin{tikzpicture}[node distance=6mm and 9mm]
  \node[term] (s) {Bắt đầu};
  \node[proc, right=of s] (p1) {Mở bài học};
  \node[dec, right=of p1] (d1) {Có bài\\tập?};
  \node[proc, below=10mm of d1] (p2) {Trả lời từng\\bài tập};
  \node[proc, left=of p2] (p3) {Chấm điểm\\từng câu};
  \node[proc, left=of p3] (p4) {Kết thúc\\lượt làm};
  \node[proc, below=9mm of p4] (p5) {Ghi\\LessonCompletion};
  \node[proc, right=of p5] (p6) {Cộng XP,\\cập nhật chuỗi};
  \node[proc, right=of p6] (p7) {Ghi điểm\\kỹ năng CEFR};
  \node[proc, below=9mm of p7] (p8) {Lập lịch ôn\\khái niệm (FSRS)};
  \node[proc, left=of p8] (p9) {Kiểm tra\\thành tích};
  \node[term, left=of p9] (e) {Kết thúc};
  \node[term, above right=4mm and 6mm of d1, fill=cSoftO, draw=cGate] (err) {Lỗi 409\\bài chưa dùng được};

  \draw[fl] (s) -- (p1);
  \draw[fl] (p1) -- (d1);
  \draw[fl] (d1) -- node[lab,near start]{không} (err);
  \draw[fl] (d1) -- node[lab,near start]{có} (p2);
  \draw[fl] (p2) -- (p3);
  \draw[fl] (p3) -- (p4);
  \draw[fl] (p4) -- (p5);
  \draw[fl] (p5) -- (p6);
  \draw[fl] (p6) -- (p7);
  \draw[fl] (p7) -- (p8);
  \draw[fl] (p8) -- (p9);
  \draw[fl] (p9) -- (e);
\end{tikzpicture}
\caption{Sơ đồ hoạt động — hoàn thành một bài học. Ba bước cuối là nơi luyện tập biến thành dữ liệu.}
\label{fig:hoatdong-baihoc}
\end{figure}

\subsection{Quy trình một lượt hội thoại với AI}

\begin{figure}[htbp]
\centering
\begin{tikzpicture}[node distance=5mm and 8mm]
  \node[term] (s) {Người học gửi câu};
  \node[proc, below=of s] (p1) {Nạp hồ sơ + lịch sử};
  \node[dec, below=of p1] (d1) {Đệm\\trúng?};
  \node[proc, right=16mm of d1] (p2) {Trả lời từ bộ đệm};
  \node[proc, below=8mm of d1] (p3) {Mở rộng đồ thị\\+ chẩn đoán lỗi (song song)};
  \node[dec, below=of p3] (d2) {Đủ\\tin cậy?};
  \node[proc, left=14mm of d2] (p4) {Hỏi làm rõ};
  \node[proc, below=8mm of d2] (p5) {Truy xuất bằng chứng\\(đồ thị + véc-tơ)};
  \node[proc, below=of p5] (p6) {Sinh phản hồi\\theo chiến lược sư phạm};
  \node[proc, below=of p6] (p7) {Ghi quan sát học tập\\vào vùng đệm};
  \node[term, below=of p7] (e) {Trả lời + điểm số};

  \draw[fl] (s) -- (p1);
  \draw[fl] (p1) -- (d1);
  \draw[fl] (d1) -- node[lab]{có} (p2);
  \draw[fl] (p2.north) |- ([yshift=3mm]p2.north) -- ++(0,0);
  \draw[fl] (d1) -- node[lab,near start]{không} (p3);
  \draw[fl] (p3) -- (d2);
  \draw[fl] (d2) -- node[lab]{$<0{,}5$} (p4);
  \draw[fl] (d2) -- node[lab,near start]{$\ge 0{,}5$} (p5);
  \draw[fl] (p5) -- (p6);
  \draw[fl] (p6) -- (p7);
  \draw[fl] (p7) -- (e);
  \draw[fl] (p4.south) |- ([yshift=-3mm]p4.south) -| (e.west);
  \draw[fl] (p2.east) -- ++(0.6,0) |- (e.east);
\end{tikzpicture}
\caption{Sơ đồ hoạt động — một lượt hội thoại. Ba nhánh kết thúc: đệm, hỏi lại, sinh mới.}
\label{fig:hoatdong-chat}
\end{figure}

\subsection{Quy trình xuất bản khoá học}

\begin{figure}[htbp]
\centering
\begin{tikzpicture}[node distance=5mm and 7mm]
  \node[term] (s) {Tải danh sách từ\\có nhãn CEFR};
  \node[proc, right=of s] (p1) {Lập chương trình\\theo bậc};
  \node[proc, right=of p1] (p2) {Sinh bài học\\và bài tập};
  \node[proc, right=of p2] (p3) {Làm sạch \&\\kiểm hình dạng};
  \node[proc, below=9mm of p3] (p4) {Lưu ở trạng thái\\bản nháp};
  \node[proc, left=of p4] (p5) {Người duyệt\\rà soát};
  \node[dec, left=of p5] (d1) {Mọi bài có\\bài tập?};
  \node[proc, left=of d1] (p6) {Sinh lại /\\sửa tay};
  \node[term, below=9mm of d1, fill=cSoftG, draw=cData] (e) {Xuất bản};

  \draw[fl] (s) -- (p1); \draw[fl] (p1) -- (p2); \draw[fl] (p2) -- (p3);
  \draw[fl] (p3) -- (p4); \draw[fl] (p4) -- (p5); \draw[fl] (p5) -- (d1);
  \draw[fl] (d1) -- node[lab]{chưa} (p6);
  \draw[fl] (d1) -- node[lab,near start]{rồi} (e);
  \draw[fl] (p6.north) |- ([yshift=4mm]p6.north) -| (p2.south);
\end{tikzpicture}
\caption{Sơ đồ hoạt động — sinh và xuất bản khoá học. Cổng chặn nằm ở bước kiểm tra bài tập.}
\label{fig:hoatdong-xuatban}
\end{figure}

\section{Quy tắc nghiệp vụ}
\label{sec:qtnv}

Mười quy tắc dưới đây là ràng buộc cứng của hệ thống. Mỗi quy tắc đều có mã để các chương
sau tham chiếu, và mỗi quy tắc đều tồn tại vì đã từng bị vi phạm trong thực tế.

\begin{longtable}{L{1.3cm} L{12.3cm}}
\toprule
\textbf{Mã} & \textbf{Nội dung}\\
\midrule
\endhead
QT-01 & Một bài học chỉ dùng được khi có ít nhất một bài tập. Khoá học chỉ xuất bản được khi mọi bài học thoả điều kiện này.\\
QT-02 & Mọi hoạt động có chấm điểm phải gửi bằng chứng tới \textbf{cả hai} động cơ: điểm kỹ năng CEFR và trạng thái khái niệm.\\
QT-03 & Tiêu thụ thụ động (nghe podcast, xem video, đọc thuần tuý) không sinh bằng chứng học tập.\\
QT-04 & Điểm kỹ năng dịch chuyển \textbf{một bước cho mỗi bài tập}, không phải một bước cho mỗi lần gọi API.\\
QT-05 & Kỹ năng của nội dung đến từ nhãn (\code{lesson.skill}, \code{course.skill}), không suy đoán từ thẻ hay tên trò chơi.\\
QT-06 & Cùng một câu trả lời không được ghi nhận hai lần: khoá chống trùng là \code{event\_id} với quan sát, và \code{source\_id} với XP.\\
QT-07 & Người học chỉ lên tối đa một bậc CEFR mỗi lần, và phải vượt cổng thi.\\
QT-08 & Thăng bậc cần cả khối lượng (số bài đã làm) lẫn độ tin cậy (mức độ hệ thống đã đo được người học).\\
QT-09 & Phản hồi lấy từ bộ đệm phải cùng bậc CEFR, cùng tiếng mẹ đẻ và cùng thế hệ trạng thái với yêu cầu hiện tại.\\
QT-10 & Quan sát học tập phải được ghi bền vững \textbf{trước} khi trả lời người dùng, không phải sau.\\
\bottomrule
\caption{Mười quy tắc nghiệp vụ ràng buộc toàn hệ thống}
\end{longtable}

\section{Ma trận nghiệp vụ và dữ liệu sinh ra}

\begin{longtable}{L{3.6cm} L{2.2cm} L{2.4cm} L{2.0cm} L{2.4cm}}
\toprule
\textbf{Nghiệp vụ} & \textbf{Điểm kỹ năng} & \textbf{Trạng thái khái niệm} & \textbf{XP} & \textbf{Thành tích}\\
\midrule
\endhead
NV1 — Bài học & Có & Có & Có & Có\\
NV2 — Ôn từ vựng & Có & Có & Có & Có\\
NV3 — Hội thoại AI & Có & Có & Có & Có\\
NV4 — Luyện nói & Có & Có & Có & Có\\
NV5 — Trò chơi & Có & Có & Có & Có\\
NV6 — Nội dung thực tế (quiz) & Có & Có & Có & Có\\
NV6 — Nội dung thực tế (nghe/đọc) & Không & Không & Không & Không\\
NV7 — Xem tiến độ & Không & Không & Không & Không\\
NV8 — Thi thăng bậc & Có & Không & Không & Có\\
\bottomrule
\caption{Nghiệp vụ nào sinh ra loại dữ liệu nào (chi tiết công thức ở Chương~9)}
\end{longtable}
